\subsection{n-th EXTERIOR POWER OF A MODULE AND ALTERNATING MULTI-LINEAR MAPPINGS}

Given two modules M, N over a commutative ring A, an alternating n-linear mapping of \(M^{n}\) into N is any n-linear mapping \(f: M^{n} \to N\) such that, for all \(p \leqslant n - 2\) ,

\[
f (u _ {1}, \dots , u _ {p}, x, x, v _ {1}, \dots , v _ {n - p - 2}) = 0\tag{8}
\]

for all \(x\) , the \(u_{i}\) ( \(1 \leqslant i \leqslant p\) ) and the \(v_{j}\) ( \(1 \leqslant j \leqslant n - p - 2\) ) in \(\mathbf{M}\) .

PROPOSITION 7. Let A be a commutative ring and M and N two A-modules. If with every A-linear mapping \(g: \bigwedge^n (\mathbf{M}) \to \mathbf{N} (n \geqslant 2)\) is associated the n-linear mapping

\[
\left(x _ {1}, x _ {2}, \dots , x _ {n}\right) \mapsto g \left(x _ {1} \wedge x _ {2} \wedge \dots \wedge x _ {n}\right)\tag{9}
\]

a bijective A-linear mapping is obtained of the A-module \(\operatorname{Hom}_{\mathbf{A}}(\bigwedge^{n}(\mathbf{M}), \mathbf{N})\) onto the A-module of alternating \(n\) -linear mappings of \(\mathbf{M}^n\) into \(\mathbf{N}\) .

We consider the canonical bijection of the A-module \(\operatorname{Hom}_{\mathbf{A}}(\mathsf{T}^{n}(\mathbf{M}),\mathbf{N})\) onto the A-module \(\mathcal{L}_{n}(\mathbf{M},\ldots,\mathbf{M};\mathbf{N})\) of all n-linear mappings of \(M^{n}\) into N, obtained by associating with every A-linear mapping \(f: T^{n}(\mathbf{M}) \to \mathbf{N}\) the n-linear mapping

\[
\bar {f} \colon (x _ {1}, \dots , x _ {n}) \mapsto f (x _ {1} \otimes x _ {2} \otimes \dots \otimes x _ {n})
\]

(II, § 3, no. 9). On the other hand, the A-linear mappings \(g \colon \bigwedge^{n}(\mathbf{M}) \to \mathbf{N}\) are in one-to-one correspondence with the A-linear mappings \(f \colon \mathsf{T}^n (\mathbf{M}) \to \mathbf{N}\) such that \(f\) is zero on \(\mathfrak{J}_n''\) , by associating with \(g\) the mapping \(f = g \circ p_n\) , where

\[
p _ {n} \colon \mathsf {T} ^ {n} (\mathbf {M}) \rightarrow \bigwedge^ {n} (\mathbf {M}) = \mathsf {T} ^ {n} (\mathbf {M}) / \mathfrak {J} _ {n} ^ {\prime \prime}
\]

is the canonical homomorphism (II, §2, no. 1, Theorem 1). But as \(\mathfrak{J}_n''\) is a linear combination of elements of the form

\[
\left(u _ {1} \otimes u _ {2} \otimes \dots \otimes u _ {p}\right) \otimes (x \otimes x) \otimes \left(v _ {1} \otimes \dots \otimes v _ {n - p - 2}\right)
\]

\((x, u_{i}, v_{j} \text{ in M}), \text{ to say that } f \text{ is of the form } g \circ p_{n} \text{ means that the corresponding } n\text{-linear function } \bar{f} \text{ satisfies (8), in other words it is alternating.}\)

The A-module \(\bigwedge^{n}(\mathbf{M})\) is called the \(n\) -th exterior power of \(\mathbf{M}\) . For every A-module homomorphism \(u: \mathbf{M} \to \mathbf{N}\) , the mapping

\[
\bigwedge^ {n} (u): \bigwedge^ {n} (\mathbf {M}) \rightarrow \bigwedge^ {n} (\mathbf {N})
\]

which coincides with \(\bigwedge(u)\) on \(\bigwedge^{n}(\mathbf{M})\) is called the \(n\) -th exterior power of \(u\) .

COROLLARY 1. For every alternating \(n\) -linear mapping \(g: \mathbf{M}^n \to \mathbf{N}\) , for every permutation \(\sigma \in \mathfrak{S}_n\) ,

\[
g \left(x _ {\sigma (1)}, x _ {\sigma (2)}, \dots , x _ {\sigma (n)}\right) = \varepsilon_ {\sigma}. g \left(x _ {1}, x _ {2}, \dots , x _ {n}\right)\tag{10}
\]

for all \(x_{i} \in \mathbf{M}\) ; moreover if \(x_{i} = x_{j}\) for two distinct indices \(i, j\) , then

\[
g (x _ {1}, x _ {2}, \dots , x _ {n}) = 0.
\]

This is an obvious consequence of Proposition 7 and no. 3, Proposition 5.

COROLLARY 2. Let \((x_{i})_{1\leqslant i\leqslant n}\) be a sequence of \(n\) elements of \(\mathbf{M}\) such that

\[
x _ {1} \wedge x _ {2} \wedge \dots \wedge x _ {n} = 0;
\]

then, for every alternating \(n\) -linear mapping \(g: \mathbf{M}^n \to \mathbf{N}\) , \(g(x_1, \ldots, x_n) = 0\) .

COROLLARY 3. Let \(f: \mathbf{M}^{n-1} \to \mathbf{A}\) be an alternating \((n-1)\) -linear form. If \((x_i)_{1 \leq i \leq n}\) is a sequence of \(n\) elements of \(\mathbf{M}\) such that \(x_1 \wedge x_2 \wedge \cdots \wedge x_n = 0\) , then

\[
\sum_ {i = 1} ^ {n} (- 1) ^ {i} f (x _ {1}, \dots , \hat {x _ {i}}, \dots , x _ {n}). x _ {i} = 0\tag{11}
\]

(where we write \(f(x_{1},\ldots ,\hat{x}_{i},\ldots ,x_{n}) = f(x_{1},\ldots ,x_{i - 1},x_{i + 1},\ldots ,x_{n})\) for \(1\leqslant i\leqslant n\) .

It suffices to prove that the n-linear mapping

\[
(x _ {1}, \dots , x _ {n}) \mapsto \sum_ {i = 1} ^ {n} (- 1) ^ {i} f (x _ {1}, \dots , \hat {x} _ {i}, \dots , x _ {n}). x _ {i}
\]

of \(M^{n}\) to M is alternating. Now, if \(x_{i} = x_{i+1}\) , all the terms in the sum on the right hand side have zero coefficients except \(x_{i}\) and \(x_{i+1}\) , since f is alternating; on the other hand, the coefficient of \(x_{i}\) is

\[
(- 1) ^ {i} f (x _ {1}, \dots , x _ {i - 1}, x _ {i + 1}, x _ {i + 2}, \dots , x _ {n})
\]

and that of \(x_{i+1}\) is \((-1)^{i+1}f(x_1, \ldots, x_i, x_{i+2}, \ldots, x_n)\) and they are inverses by hypothesis.

Remark. An element \(z\) of \(\mathsf{T}^n (\mathbf{M})\) is called a (contravariant) skew-symmetric tensor of order \(n\) if \(\sigma .z = \varepsilon_{\sigma}z\) for every permutation \(\sigma \in \mathfrak{S}_n\) (cf.~§6, no. 3, Remark); these elements form a sub-A-module \(\mathsf{A}_n^\prime (\mathbf{M})\) of \(\mathsf{T}^n (\mathbf{M})\) . For all \(z\in \mathsf{T}^n (\mathbf{M})\) , we write \(\pmb {a}.z = \sum_{\sigma \in \mathfrak{S}_n}\varepsilon_\sigma (\sigma .z)\) and call \(\pmb {a}.z\) the skew-symmetrization of \(z\) ; as \(\varepsilon_{\sigma \tau} = \varepsilon_{\sigma}\varepsilon_{\tau}\) , it is immediately seen that \(\pmb {a}.z\) is a skew-symmetric tensor and therefore \(z\mapsto \pmb {a}.z\) is an endomorphism of \(\mathsf{T}^n (\mathbf{M})\) whose image \(\mathsf{A}_n''(\mathbf{M})\) is contained in \(\mathsf{A}_n^\prime (\mathbf{M})\) ; in general \(\mathsf{A}_n''(\mathbf{M})\neq \mathsf{A}_n^\prime (\mathbf{M})\) (Exercise 8). If \(z\) is a skew-symmetric tensor, then \(\pmb {a}.z = n!z\) ; hence, when \(n!\) is invertible in A, the endomorphism \(z\mapsto (n!)^{-1}\pmb {a}.z\) is a projector of \(\mathsf{T}^n (\mathbf{M})\) whose image is

\[
\mathbf {A} _ {n} ^ {\prime} (\mathbf {M}) = \mathbf {A} _ {n} ^ {\prime \prime} (\mathbf {M}).
\]

Moreover the kernel of this projector is just \(\mathfrak{J}_n''\) ; for attention may obviously be confined to the case where \(n \geqslant 2\) , hence 2 (a divisor of \(n!\) ) is invertible in A and \(x \otimes x = 2^{-1}(x \otimes x + x \otimes x)\) ; therefore \(\mathfrak{J}_n''\) is generated by the elements \(z + \rho . z\) , where \(\rho\) is a transposition exchanging two consecutive numbers in \([1, n]\) ; moreover, for two permutations \(\sigma, \tau\) in \(\mathfrak{S}_n\) , we can write

\[
z - \varepsilon_ {\sigma \tau} ((\sigma \tau). z) = z - \varepsilon_ {\tau} (\tau . z) + \varepsilon_ {\tau} (\tau . z - \varepsilon_ {\sigma} \sigma . (\tau . z))
\]

whence it follows that \(z - \varepsilon_{\sigma}(\sigma . z) \in \mathfrak{J}_{n}^{\prime \prime}\) for all \(z \in \mathbf{T}^{n}(\mathbf{M})\) and \(\sigma \in \mathfrak{S}_{n}\) . Therefore (always assuming \(n!\) is invertible in A), it is seen that

\[
z - (n!) ^ {- 1} \boldsymbol {a}. z = \sum_ {\sigma \in \mathfrak {S} _ {n}} (n!) ^ {- 1} (z - \varepsilon_ {\sigma} (\sigma . z)) \in \mathfrak {J} _ {n} ^ {\prime \prime}
\]

for all \(z \in \mathbf{T}^n(\mathbf{M})\) , which establishes our assertion.

When \(n!\) is invertible in A, the submodules \(\mathsf{A}_n'(M)\) and \(\mathfrak{J}_n''\) of \(T^n(M)\) are therefore supplementary and the restriction to \(\mathsf{A}_n'(M)\) of the canonical homomorphism \(T^n(M) \to \bigwedge^n(M) = T^n(M)/\mathfrak{J}_n''\) is an A-module isomorphism, which allows us in the case in question to identify skew-symmetric tensors of order \(n\) with the elements of the \(n\) -th exterior power of M. Note here also that this identification is not compatible with multiplication, the product in \(T(M)\) of two skew-symmetric tensors not being skew-symmetric in general.

